\chapter*{Abstract}
\addcontentsline{toc}{chapter}{Abstract}

Early-stage portfolio reporting combines company submissions with public evidence that may be incomplete, delayed or inconsistent. The current process has three broad stages: company information is collected, aggregated into an executive summary, and prepared for review before use. The main risks are incorrect company identity, confusion between blank and zero values, evidence from the wrong period, unresolved conflicts and claims without adequate support. This dissertation evaluates an implemented evidence-first workflow that addresses these risks through explicit data rules, traceable source records and a separately recorded verification decision. It asks how the workflow preserves identity, reporting time, missing states, sources and approval boundaries (RQ1), and what changes when a separate verification stage controls which claims enter the report (RQ2).

The prototype is a local reporting application that links each claim to its source, period, missing-data reason and approval. It separates planning, identity checks, collection, extraction, normalisation, verification, composition and approval into recorded roles with limited authority. Recorded hand-offs show who or what produced, checked and approved each statement. The roles do not imply that several agents are generally better than a simpler system. A controlled public-web route was also tested with fake providers and local fixtures, but not against a live company or website.

The evaluation used fourteen labelled fictional cases and repeated the comparison three times. Without the separate verifier, the baseline produced eleven claims: five were supported and six were not. With verification, the workflow produced the five supported claims and held all six unsupported candidates. In metric terms, precision (the share of produced claims that were supported) increased from 0.455 to 1.000, while recall (the share of required supported claims retained) remained 1.000. The combined F1 measure increased from 0.625 to 1.000, source-record completeness increased from 0.818 to 1.000, and the unsupported-claim rate fell from 0.545 to 0.000. The verifier retained an observed zero with matching evidence and held contradictory, stale and untrusted cases. These results establish the programmed mechanism on designed cases.

The contribution is therefore an inspectable evidence-first implementation together with the observed effect of its verification gate on designed cases. Because the cases and the rules were written together, the result shows that a programmed gate behaves as specified and not that the system performs well on real companies, saves analyst time or is ready for production. Manual reporting, participant review, live-web evaluation and a managed-platform comparison all lacked the observations they required, so they are set out as a prospective pilot rather than reported as findings.
